Three GLP-1 based medications are FDA-approved for chronic weight management in adults with obesity (BMI $\ge$ 30) or overweight (BMI $\ge$ 27) with at least one weight-related chronic condition. Saxenda (liraglutide), a GLP-1 receptor agonist, was approved on December 23, 2014. Wegovy (semaglutide), also a GLP-1 receptor agonist, was approved on June 4, 2021. The newest drug is a dual GIP/GLP-1 receptor agonist called Zepbound (tirzepatide), approved on November 8, 2023. Saxenda is administered as a once-daily injection; Wegovy and Zepbound are typically given as once-weekly injections.\footnote{In December 2025, the FDA approved the Wegovy Pill, a once-daily oral form of semaglutide; it became available in pharmacies in January 2026. } All three drugs start with a low dose and are titrated up to a maintenance dose to mitigate side effects, which include nausea, vomiting, diarrhea/constipation, abdominal pain, and decreased appetite. An umbrella review of meta-analyses suggests side effects remain common and are an important factor in discontinuation of GLP-1 therapy \citep{kong2026glp1umbrella}. Weight loss efficacy has improved with each generation: in major trials, percent weight loss relative to placebo was 5.4 percentage points lower for Saxenda patients \citep{pi2015randomized}, 12.5 percentage points lower for Wegovy patients \citep{Wilding2021}, and 17.8 percentage points lower for 15mg Zepbound patients \citep{Jastreboff2022}.

Prescription drug pricing in the United States is often opaque. Manufacturers, pharmacy benefit managers, and insurers use discounts, rebates, and cross-subsidies that obscure the true cost of a drug to any given payer. With that caveat, Table~\ref{tab:consumer_prices} shows three publicly available measures of the price of obesity-indicated GLP-1 medications. The National Average Drug Acquisition Cost (NADAC) is based on a CMS survey of pharmacy reports of purchase invoices that reflects what pharmacies pay wholesalers for each drug. The Federal Supply Schedule (FSS) price is a contract price negotiated between the VA and drug manufacturers for eligible federal purchasers. The Big~4 price is a further statutory discount available to the VA, Department of Defense, Coast Guard, and Public Health Service---typically the lowest price available to any government purchaser. At NADAC prices, monthly costs are approximately \$1,300 for Saxenda and Wegovy and \$1,050 for Zepbound. These prices do not include coupons or rebates, which may be substantial. But even at the Big~4 price, monthly costs are about \$1,000 for Saxenda and Wegovy and  \$800 for Zepbound. And these prices have been relatively stable since each drug's launch. Although the nominal prices are high, GLP-1 drugs represent a major improvement in the quality of medical care for obesity. Quality improvement across the three generations of drugs is substantial, even in more recent formulations. The monthly NADAC divided by the percentage point treatment effect---a crude quality-adjustment---implies that the price of a one percent reduction in body weight fell from approximately \$241 with Saxenda to \$104 with Wegovy, and to only \$59 with Zepbound. Nevertheless, at current prices, the absolute cost of covering these drugs for the Medicaid population is apt to be a substantial new fiscal commitment.

Under the Medicaid Drug Rebate Program (MDRP), in order for \emph{any and all of their products} to be covered by Medicaid, pharmaceutical manufacturers must enter into rebate agreements with the federal government and pay quarterly rebates to states according to a formula that depends on confidential information about their average and best prices and product-specific inflation. In exchange, states participating in the MDRP are required to cover all FDA-approved drugs made by participating manufacturers. However, Section~1927(d)(2) of the Social Security Act carves out several categories of drugs that states \emph{may} exclude from coverage, including ``agents when used for anorexia, weight loss, or weight gain'' \citep{SSA1927}. This statutory exclusion means that obesity-indicated GLP-1 medications are not automatically covered by Medicaid. States must make an affirmative decision to add them to their formularies. The result is substantial cross-state variation in whether Medicaid pays for GLP-1 obesity drugs, which forms the basis of our identification strategy.

We compiled data on Medicaid coverage for obesity-indicated GLP-1 medications, searching over years 2019 to 2026, tracking both the timing of coverage and the use of utilization management strategies for each drug in each state.\footnote{Our work builds on earlier work by the Kaiser Family Foundation, which cataloged coverage of obesity-indicated GLP-1 medications in mid-2024 and 2026 \citep{kff2024glp1,williams2026KFF}.} Our primary source is the Medicaid Preferred Drug List (PDL) for each state and year. PDLs designate preferred drugs accessible without prior authorization, and typically specify the conditions---such as prior authorization, step therapy, or quantity limits---under which non-preferred drugs may be obtained \citep{simon_etal_2009, dolan_tian_2019}. We reviewed PDLs for all 50 states from 2019 through 2026, recording whether Wegovy, Saxenda, or Zepbound appeared with an explicit obesity indication and under what utilization management conditions. When PDL listings were ambiguous, we corroborated with secondary sources including pharmacy program amendments, Pharmacy and Therapeutics Committee minutes, prior authorization criteria documents, and state Medicaid pharmacy portal NDC lookups. The replication archive includes archived PDLs, supplementary policy documents, a state-specific SDUD utilization graph, and a summary memorandum with primary source citations.\footnote{A browsable page with state-by-state evidence sources and full citations is available at \url{https://gist.github.com/coadywing/14061c2322ad07fcd042d1aca7a65694}.}

Our coverage dates reflect fee-for-service (FFS) Medicaid policy. While most Medicaid beneficiaries under 65 are enrolled in managed care organizations (MCOs), FFS PDL policy is an increasingly reliable signal of MCO coverage. States have rapidly adopted uniform PDL requirements that bind MCOs to the same formulary as FFS: as of July 2019, 16 MCO states had uniform PDLs, rising to 19 of 30 MCO pharmacy carve-in states by July 2023 \citep{gifford_etal_2020, healthmanagement_2024}. Even in states without uniform PDLs, most require MCO prior authorization to be no more restrictive than FFS \citep{gifford_etal_2020}. We validated the classifications by examining FFS and MCO prescription volume in the State Drug Utilization Data (SDUD), confirming that the first quarter with substantial utilization matched or immediately followed our PDL-documented coverage date for both FFS and MCO utilization.

Michigan was the first state to cover GLP-1 obesity medications, adding Saxenda and Wegovy to its Medicaid formulary in February 2022 (2022-Q1). New Hampshire followed in September 2022 (2022-Q3), and Rhode Island in October 2022 (2022-Q4). Four states---California, Delaware, Pennsylvania, and Virginia---adopted coverage in January 2023 (2023-Q1). By the end of 2022, three states had coverage; by 2023-Q1, seven; by 2023-Q4, ten; and by mid-2025, 17 states had adopted coverage for at least one GLP-1 obesity drug. Early adopters typically covered Saxenda and Wegovy simultaneously. Zepbound, approved in November 2023, was subsequently added by covering states beginning in early 2024; Pennsylvania and Virginia were among the first to add Zepbound in January 2024. Figure~\ref{fig:coverage_map} displays the geographic distribution of coverage adoption. Table~\ref{tab:design} summarizes the timing of first coverage across states, and Table~\ref{tab:coverage_details} provides the complete set of coverage start dates by state and drug.

Of the 17 states that covered at least one GLP-1 obesity drug during our study period, 16 implemented prior authorization (PA) at the time of initial coverage; Rhode Island was the exception, initially covering Wegovy and Saxenda without PA in October 2022 before adding PA requirements in 2025. California employed quantity limits rather than prior authorization. Prior authorization criteria vary but typically require a documented obesity diagnosis (BMI $\ge$ 30, or $\ge$ 27 with comorbidities) and evidence of prior lifestyle modification. Table~\ref{tab:pa_utilization} classifies covering states into tiers by the stringency of their PA requirements. Three states (Minnesota, Pennsylvania, and Tennessee) implemented both prior authorization and quantity limits from the start of coverage.

Very recently, four states have rescinded coverage for obesity-indicated GLP-1 medication. California, New Hampshire, Pennsylvania, and South Carolina announced coverage terminations effective in late 2025 and early 2026, citing budgetary pressures. Two states have contined with coverage but have adopted more restrictive eligibility criteria. Massachusetts restricted eligibility to individuals with BMI of 35 or greater, and Virginia restricted eligibility to individuals with BMI of 40 or greater. These changes do not affect the results of our analysis, which analyzes outcomes out to 2025 Q2.

